\documentclass[11pt,a4paper]{article}
\usepackage[T1]{fontenc}
\usepackage[utf8]{inputenc}
\usepackage{amsmath,amssymb,amsthm}
\usepackage[margin=1in]{geometry}
\usepackage[colorlinks=true,linkcolor=blue,urlcolor=blue]{hyperref}
\theoremstyle{plain}
\newtheorem{theorem}{Theorem}
\newtheorem{lemma}[theorem]{Lemma}
\newtheorem{proposition}[theorem]{Proposition}
\newtheorem{corollary}[theorem]{Corollary}
\theoremstyle{definition}
\newtheorem{remark}[theorem]{Remark}

\title{The empirical recurrence for OEIS A232593, proved by transfer matrix}
\author{Adrian Perez Fontelles\\ \small Independent researcher}
\date{4 September 2026}
\begin{document}
\maketitle

\begin{abstract}
OEIS A232593 counts arrays of width $5$ over $\{0,1,2\}$ in which every cell sees its own value plus and minus one among its neighbours. The entry carries an
empirical recurrence of order $49$ contributed by R.~H.~Hardin. It is true, and it is
decidable rather than empirical: what a cell asks reaches one row above it and one row below, so the
condition on a whole row is settled by $3$ consecutive rows and the count is a walk
count on $S=1234$ states.
\end{abstract}

\noindent\small 2020 Mathematics Subject Classification. 05A15, 05B20, 15A18.\normalsize

\section{The conjecture}

OEIS A232593 is ``Number of (4+1)X(n+1) 0..2 arrays with every element next to itself plus and minus one within the range 0..2 horizontally or antidiagonally, with no adjacent elements equal.''. It has offset $1$ and begins
\[
6, 456, 122, 6406, 3002, 93236, 68072, 1364998,\ \dots
\]
The entry states, as a conjecture contributed by its author and never marked settled:
\begin{quote}
Empirical: a(n) = 20*a(n-2) +16*a(n-3) -92*a(n-4) -260*a(n-5) +82*a(n-6) +1038*a(n-7) +1278*a(n-8) -1350*a(n-9) -4927*a(n-10) -3536*a(n-11) +5500*a(n-12) +12930*a(n-13) +8739*a(n-14) -8650*a(n-15) -22849*a(n-16) -24728*a(n-17) -3836*a(n-18) +33686*a(n-19) +62865*a(n-20) +36392*a(n-21) -44749*a(n-22) -110478*a(n-23) -67221*a(n-24) +47184*a(n-25) +126718*a(n-26) +76044*a(n-27) -30029*a(n-28) -98898*a(n-29) -68985*a(n-30) +2372*a(n-31) +49486*a(n-32) +46722*a(n-33) +19080*a(n-34) -1646*a(n-35) -13297*a(n-36) -15236*a(n-37) -12882*a(n-38) -4684*a(n-39) +949*a(n-40) +3512*a(n-41) +2385*a(n-42) +1194*a(n-43) +233*a(n-44) -6*a(n-45) -67*a(n-46) -18*a(n-47) -2*a(n-48) +2*a(n-49) for n>54
\end{quote}
As of the ``Last modified'' line on the live entry (Jul 23 2025 07:31:32, revision 6) this is still
recorded as empirical, and nothing on the entry records it as proved.

\section{What is being counted}

The arrays have $5$ columns and a growing number of rows, with entries in $\{0,1,2\}$. Fix the
offsets the entry names,
\[
\Delta=\{\,(-1,0), (-1,1), (1,-1), (1,0)\,\},
\]
and for a cell $(i,j)$ write
\[
S(i,j)=\{\,(i+\delta_1,\,j+\delta_2) : (\delta_1,\delta_2)\in\Delta\,\}\cap
\text{(the cells of the array)}
\]
for the neighbours it actually has. Only offsets landing inside the array count, so a cell on
the border simply has fewer neighbours. The entry writes the array with $n$ as the number of COLUMNS; transposing it exchanges the two coordinates of every offset, and the offsets above are written in the transposed frame.

The condition is that every cell has, among its neighbours, each of its own value plus one and its own value minus one that still lies in the range:
\[
t\in\{x(i,j)-1,\ x(i,j)+1\}\cap\{0,\dots,2\}\ \Longrightarrow\ t\in\{x(a,b):(a,b)\in S(i,j)\}.
\]
At the ends of the range only one of the two values exists and only that one is demanded.

The entry adds that no two cells at one of these offsets are equal.

\section{A window of 3 rows}

The offsets reach one row up and one row down, so everything the condition asks of one row is
decided by that row together with one row above it and one row below --- a window of
$3$ consecutive rows, and nothing wider. Take the array one column at a time and
let the state be the last 2 of them, with a distinguished start standing for the array before its
first row, so that a walk of $n$ steps is an array of $n$ rows.

Appending a row completes exactly one window, and the row it settles is the one just above the bottom of that window; a step is
allowed when that row passes. Each row of the finished array is therefore tested exactly once.

When the array stops, the last row of it has never been tested, because the
conditions reach one row below. So the end vector is NOT all ones: a state is accepted exactly when
those rows pass with nothing below them, which is decided by running the same test on the window
with absent rows appended --- the same rule that governs a cell at the bottom edge in
Section 2.
Thus
\[
a(n)\;=\;\iota^{\!\top}M^{\,j}\tau ,\qquad S=1234,
\]
for the transfer matrix $M$ on those $S$ states. The walk index $j$ and the entry's own $n$
differ by a constant, read off by matching the published terms rather than assumed. In
particular $a$ is $C$-finite.

\section{The proof}

\begin{lemma}\label{lem:crit}
Let $q(t)=t^{49}-\sum_i c_it^{49-i}$ be the characteristic polynomial of the
conjectured recurrence. The residual $a(n)-\sum_ic_ia(n-i)$ equals
$\iota^{\!\top}M^{\,j}q(M)\tau$ for the corresponding walk index $j$, so the recurrence
holds from some point on if and only if $\iota^{\!\top}M^{j}q(M)\tau=0$ for all large $j$,
and it is enough to examine $j<S$.
\end{lemma}

\begin{proof}
Factor $M^{j}$ out of $M^{j+49}-\sum_ic_iM^{j+49-i}$. For the bound, $M^{S}$
is an integer combination of $I,M,\dots,M^{S-1}$ by Cayley--Hamilton, so the numbers
$u_j=\iota^{\!\top}M^{j}q(M)\tau$ obey the monic recurrence given by the characteristic
polynomial of $M$; $S$ consecutive zeros force every later one, and the last nonzero $u_j$
pins the threshold exactly.
\end{proof}

\begin{theorem}
\[
a(n)\;=\;20\,a(n-2) + 16\,a(n-3) - 92\,a(n-4) - 260\,a(n-5) + 82\,a(n-6) + 1038\,a(n-7) + 1278\,a(n-8) - 1350\,a(n-9) - 4927\,a(n-10) - 3536\,a(n-11) + 5500\,a(n-12) + 12930\,a(n-13) + 8739\,a(n-14) - 8650\,a(n-15) - 22849\,a(n-16) - 24728\,a(n-17) - 3836\,a(n-18) + 33686\,a(n-19) + 62865\,a(n-20) + 36392\,a(n-21) - 44749\,a(n-22) - 110478\,a(n-23) - 67221\,a(n-24) + 47184\,a(n-25) + 126718\,a(n-26) + 76044\,a(n-27) - 30029\,a(n-28) - 98898\,a(n-29) - 68985\,a(n-30) + 2372\,a(n-31) + 49486\,a(n-32) + 46722\,a(n-33) + 19080\,a(n-34) - 1646\,a(n-35) - 13297\,a(n-36) - 15236\,a(n-37) - 12882\,a(n-38) - 4684\,a(n-39) + 949\,a(n-40) + 3512\,a(n-41) + 2385\,a(n-42) + 1194\,a(n-43) + 233\,a(n-44) - 6\,a(n-45) - 67\,a(n-46) - 18\,a(n-47) - 2\,a(n-48) + 2\,a(n-49)
\]
for every $n>54$.
\end{theorem}

\begin{proof}
The vector $w=q(M)\tau$ was formed by $49$ matrix--vector products in exact integer
arithmetic and $\iota^{\!\top}M^{j}w$ evaluated for $j=0,1,\dots$, again exactly; those
integers vanish from the index corresponding to $n=54+1$ onwards and the run continued
until the working vector was identically zero, which settles every larger $j$ at once.
\end{proof}

\section{Verification}

Four checks, each independent of the algebra above.

First, the model reproduces all $21$ terms the entry publishes, exactly, in integer
arithmetic.

Second, the count was recomputed for the smallest arrays straight from the definition: every
array of width $5$ over $\{0,1,2\}$ with a few rows was written out cell by cell and each cell
tested against the neighbours it has. No window, no state and no end vector are involved in that
computation, and the published terms came back.

Third, the conjectured recurrence was evaluated directly on the published terms, with no
matrices involved, and holds wherever the proved range and the data overlap.

Fourth, the annihilation test of Lemma~\ref{lem:crit} was carried out in exact integer
arithmetic on the whole vector rather than on a sampled prefix.

\begin{thebibliography}{9}
\bibitem{oeis} The OEIS Foundation, \emph{The On-Line Encyclopedia of Integer
Sequences}, \url{https://oeis.org}, 2026. Sequence A232593.
\bibitem{stanley} R.~P.~Stanley, \emph{Enumerative Combinatorics}, Volume 1, 2nd ed.,
Cambridge University Press, 2011. (Transfer-matrix method, Section 4.7.)
\bibitem{horn} R.~A.~Horn and C.~R.~Johnson, \emph{Matrix Analysis}, 2nd ed., Cambridge
University Press, 2013.
\end{thebibliography}

\end{document}
